% EcoMD @ NeurIPS Workshop on Generative AI in Finance (non-archival, <=4pp main body).
% Build: latexmk -pdf main.tex   (or pdflatex x2 + bibtex).
% Format: drop the official neurips_2024.sty/.bst into this dir for the submission format;
% otherwise this file compiles with a NeurIPS-approximating fallback (see preamble).
\documentclass{article}

% --- NeurIPS format if available, else a close fallback so it always compiles ---
\IfFileExists{neurips_2024.sty}{%
  \usepackage{neurips_2024}            % submission (anonymized + line numbers); add [final] for camera-ready
}{%
  \usepackage[letterpaper,margin=1in]{geometry}
  \usepackage{times}
  \usepackage{titlesec}
  \titleformat{\section}{\normalfont\large\bfseries}{\thesection}{0.6em}{}
  \titleformat{\subsection}{\normalfont\normalsize\bfseries}{\thesubsection}{0.6em}{}
  \setlength{\parskip}{2pt}
}

\usepackage[utf8]{inputenc}
\usepackage[T1]{fontenc}
\usepackage{amsmath,amssymb}
\usepackage{graphicx}
\usepackage{booktabs}
\usepackage{xcolor}
\usepackage{microtype}
\usepackage[numbers,sort&compress]{natbib}
\usepackage{hyperref}
\hypersetup{colorlinks=true,linkcolor=blue,citecolor=blue,urlcolor=blue}

% figure placeholder that compiles without the PDF; swap real figures into ./figures/
\newcommand{\figbox}[2]{%
  \IfFileExists{#1}{\includegraphics[width=\linewidth]{#1}}{%
  \fbox{\parbox[c][3.2cm][c]{0.96\linewidth}{\centering\small\itshape #2}}}}

\title{EcoMD: A Differentiable Generative Market Simulator,\\ and a Burn-In Pitfall in Scoring It}
\author{Anonymous Submission}
\date{}

\begin{document}
\maketitle

\begin{abstract}
Generative simulators of financial markets are used for scenario generation and stress testing, and
are evaluated by how well their rollouts reproduce stylized facts such as heavy-tailed returns and
volatility clustering. We present \textbf{EcoMD}, a \emph{differentiable} particle simulator whose
differentiability enables gradient-based calibration and a \emph{controllable shock interface} for
scenario generation, and we report a \textbf{measurement pitfall in how such generators are
evaluated}. Standard practice scores the full rollout without discarding an initial warm-up; we show
this misreads a startup \emph{equilibration transient} as a stationary fat-tail match, inflating the
apparent tail fidelity (the Hill tail index rises from $3.9$ to $6.6$ once warm-up is discarded; the
steady state is light-tailed). We give the corrected protocol. Using the shock interface, EcoMD
generates \emph{coherent-liquidation} stress scenarios with a distinctive order-flow signature
(order-flow-imbalance memory jumps from $0.02$ to $0.99$ at the shock and relaxes), whereas exogenous
price gaps are inert---a controllability result. Finally, we bound fidelity honestly: the model's
heavy tails are \emph{transient}, while real one-minute crypto crash tails are \emph{stationary}
(pre-registered null test over five episodes). EcoMD is a useful controllable generator, but its tail
realism must be validated against real stationarity before use.
\end{abstract}

\section{Introduction}
Generative models of markets---agent-based simulators \citep{byrd2020abides}, neural stochastic
differential equations, and GANs over price paths---are increasingly used for scenario generation,
stress testing, and data augmentation. Their quality is judged almost universally by
\emph{stylized-fact matching}: a rollout is scored on whether its returns reproduce empirical
regularities such as the heavy (inverse-cubic) tail and volatility clustering \citep{cont2001,gabaix2003}.
The \emph{scoring procedure itself}, however, is rarely audited: whether a high stylized-fact score
certifies a \emph{stationary} property of the generator or an artifact of how the rollout is measured
is seldom checked. This motivates two practical questions: is the evaluation rigorous, and how
\emph{controllable} are the generated scenarios---which interventions produce which market regimes?

We study both through \textbf{EcoMD}, a differentiable particle simulator in which agents evolve under
overdamped Langevin dynamics with learned pairwise potentials, and prices form from aggregate excess
demand. Differentiable agent-based models are not new \citep{bouchaud1998,dyer2024,chopra2023}; we use
EcoMD's differentiability for a controllable shock interface and report three findings about generating
and evaluating market scenarios, the first of which is largely model-agnostic:
\begin{itemize}\setlength{\itemsep}{1pt}
  \item \textbf{An evaluation pitfall and its fix} (\S\ref{sec:pitfall}): warm-up-inclusive rollout
    scoring inflates fat-tail fidelity by reading an equilibration transient as a stationary law.
  \item \textbf{Controllability} (\S\ref{sec:control}): a coherent-liquidation shock drives a
    distinctive, relaxing order-flow scenario; an exogenous price gap is inert.
  \item \textbf{Honest fidelity bounds} (\S\ref{sec:fidelity}): generated tails are transient while
    real crash tails are stationary---a scenario-realism caveat, established with a pre-registered test.
\end{itemize}

\section{EcoMD: a differentiable generative simulator}\label{sec:ecomd}
EcoMD evolves $N$ agents $s_i \in \mathbb{R}^d$ by overdamped Langevin dynamics,
$\dot{s}_i = -\nabla_{s_i} U(\{s\}) + \sqrt{2T}\,\xi_i$, with a learned interaction potential $U$
(an equivariant message-passing form in the spirit of MACE \citep{batatia2022}) and temperature $T$.
The log-price increment is driven by aggregate excess demand,
$\Delta p_t = \beta\, \mathrm{ED}_t + \sigma\eta_t$ with $\mathrm{ED}_t = \kappa\sum_i \Delta s_{i,0}$,
i.e. the net signed agent flow. The model is trained to match a panel of stylized facts and is fully
differentiable, which buys two things a non-differentiable agent-based simulator
\citep{byrd2020abides} does not have: gradient-based calibration, and a clean \emph{shock interface}
for controllable interventions (\S\ref{sec:control}). Figure~\ref{fig:pitfall} (left) shows sample
paths and the stylized facts EcoMD reproduces. The rest of the paper concerns how to evaluate and
steer this generator.

\section{The evaluation pitfall}\label{sec:pitfall}
\textbf{Setup.} Rollout-based generators are scored by computing stylized facts on a generated series
of length $T$ and comparing to empirical targets. The fat-tail fact is the Hill tail index
$\alpha$ of returns \citep{hill1975}: $\alpha\!\approx\!3$ is the empirical inverse-cubic regime,
lower is heavier. The standard recipe computes $\alpha$ on the \emph{entire} rollout.

\textbf{The artifact.} A freshly initialized rollout is out of equilibrium; it relaxes to its steady
state over an initial transient. During that transient the aggregate flow is heavy-tailed, so a
warm-up-inclusive Hill estimate reports a heavy tail---and the scorer reads this as a successful match
to the empirical cube law. But it is a startup artifact. Discarding the first $\sim\!50$ steps moves
the estimate sharply toward the light steady state: on EcoMD, the baseline Hill index rises from
$3.87$ to $6.55$, and a concave-impact variant tuned to ``match'' the cube law rises from $5.05$ to
$9.26$ (Table~\ref{tab:warmup}; Fig.~\ref{fig:pitfall}, right). The steady state is \emph{light}-tailed
($\alpha\!\approx\!4.7$ for excess demand); the apparent fat-tail match was burn-in inflation.
Crucially the artifact is \emph{rewarded}, not caught: the warm-up-inclusive (heavier) reading sits
closer to the empirical cube-law target than the true light steady state, so a stylized-fact scorer
registers a \emph{better} match---the practitioner sees success, not a bug. Its low cross-seed
variance reinforces the illusion of a robust law, which is why the pitfall survives standard practice.

\begin{table}[t]
\centering
\small
\caption{Hill tail index with vs.\ without a warm-up discard (EcoMD, $T\!=\!8000$). A higher value is
a \emph{lighter} tail. Warm-up-inclusive scoring reads the equilibration transient as a stationary
cube-law match; discarding warm-up reveals a light steady state.}
\label{tab:warmup}
\begin{tabular}{lccc}
\toprule
model & no discard & discard 50 & $\Delta$ \\
\midrule
baseline & 3.87 & 6.55 & $+2.68$ \\
concave-impact (tuned to ``solve'') & 5.05 & 9.26 & $+4.21$ \\
\bottomrule
\end{tabular}
\end{table}

\textbf{The fix.} We recommend a simple protocol: (i) compute each stylized fact over a sliding window
and plot it against warm-up-discard length $W$; (ii) discard up to the plateau $W^\star$ where the
estimate stabilizes; (iii) report the score at $W^\star$ \emph{with} the sensitivity curve, and flag
any fact that shifts by more than a set tolerance over $[0,W^\star]$ as transient-contaminated. This is
a few-line change to any rollout-scoring pipeline, yet it changes qualitative conclusions: under it,
EcoMD's earlier ``cube-law reproduction'' is withdrawn. We recommend it as default hygiene for
rollout-based market generators.

\begin{figure}[t]
\centering
\figbox{figures/fig1_pitfall.pdf}{Fig.~1 placeholder --- (left) EcoMD sample paths + stylized-fact
panel and the shock-interface schematic; (right) Hill tail index vs.\ warm-up discard length:
warm-up-inclusive scoring (left edge) reads a heavy ``cube-law'' tail that vanishes once the
equilibration transient is dropped.}
\caption{EcoMD as a generator, and the warm-up-scoring pitfall. The headline of this paper:
fat-tail ``fidelity'' measured on full rollouts is inflated by a startup equilibration transient.}
\label{fig:pitfall}
\end{figure}

\section{Controllable stress scenarios}\label{sec:control}
The shock interface injects a scheduled intervention mid-rollout, turning EcoMD into a controllable
scenario generator. We contrast two interventions.

\textbf{Coherent liquidation} (a panic / forced-selling analogue) displaces a fraction of agents'
latent positions in a common direction. This drives a clean, \emph{relaxing} stress scenario: the
heavy transient tail re-appears ($\alpha$ for excess demand drops from $4.7$ to $\approx\!0.5$ and
recovers within $\sim\!10^3$ steps), with a graded, sigmoid dose-response that generalizes across five
assets (equities, gold, crypto, FX). It also leaves a distinctive \emph{order-flow} signature: the
lag-1 autocorrelation of order-flow imbalance jumps from $0.02$ to $0.99$ at the shock---order flow
becomes transiently coherent and persistent---and relaxes back, while the imbalance magnitude and
its extreme-coordination saturation burst and decay in step (Fig.~\ref{fig:control}, left).

\textbf{Exogenous price gap.} Injecting an exogenous return into the price (a news-shock analogue) is,
by contrast, \emph{inert}: across magnitudes up to $12\sigma$ the excess-demand tail stays light
($\alpha\!\approx\!4.6$), the return tail only nudges (from $\approx\!7.7$ to $6.0$, still far from
heavy), and the order-flow imbalance is statistically identical to the unshocked control. A price
gap does not propagate into a coordinated order-flow response in this model.

The initial coherence is partly imposed by the intervention; the non-trivial, emergent content is its
\emph{finite-time relaxation} and the sharp contrast with the inert price gap---an exogenous price
move, the obvious ``stress'' lever, does nothing here. The practical takeaway is that
\emph{controllability is intervention-specific}: the lever that produces realistic stress is
coordinated participant behavior, not a price shock. The coherent-liquidation order-flow signature is
also a concrete, falsifiable prediction to test against real crash order-flow data, complementing
fluctuation-theorem analyses of market cascades \citep{maskawa2025}.

\section{Honest fidelity bounds}\label{sec:fidelity}
A generator is only as useful as its fidelity is understood. EcoMD's heavy tails are a \emph{driven
transient}; are real market tails the same? We test this directly. On five real one-minute crypto
crash episodes (COVID-2020, the May-2021 selloff, the June-2022 deleveraging, the Terra/Luna
collapse, and FTX), we pre-registered a statistic $\Delta\alpha = \alpha(\text{crash})-\alpha(\text{pre})$
on volatility-standardized returns and compared it to a null distribution built from a long calm
window (Fig.~\ref{fig:control}, right). If crashes drove a heavier tail, $\Delta\alpha$ would be
strongly negative. Instead the pooled effect is $z=+1.03$ (slightly \emph{lighter}, opposite to the
prediction); only one of five episodes shows a significant heavier tail, within the false-positive
rate for five tests. Real return tails are \emph{stationary}---approximately cube-law in calm and
crash alike, consistent with the inverse-cubic universality \citep{gabaix2003,plerou1999,toth2018}.
This uses the five episodes available with free intraday data (crypto); a broader equity test awaits
paid high-frequency data, but the result is robust to volatility standardization and already matches
the established stationarity of the return tail.

Hence EcoMD generates a heavy tail that is a \emph{non-equilibrium} response, not the real
\emph{stationary} heavy tail. For scenario generation this is a sharp caveat: the simulator reproduces
fat tails only transiently and under coordinated driving, so its generated tails should not be treated
as realistic stationary risk without external validation. It also localizes what the model class is
missing---a stationary heavy-tail mechanism---which we leave to future work (a natural candidate is a
heterogeneous, heavy order-flow source).

\begin{figure}[t]
\centering
\figbox{figures/fig2_control_fidelity.pdf}{Fig.~2 placeholder --- (left) controllable stress scenario:
order-flow-imbalance memory/|imbalance|/saturation burst-and-relax under coherent liquidation vs.\
flat under an exogenous price gap; (right) the real-data fidelity test: per-episode and pooled
$\Delta\alpha$ against the calm null --- real crash tails do not heavy-up.}
\caption{Controllability (left) and the transient-vs-stationary fidelity boundary (right).}
\label{fig:control}
\end{figure}

\section{Discussion}
We offer three takeaways for generative AI in finance. (i) \emph{Evaluate with warm-up hygiene}:
rollout-based fat-tail scores are inflated by startup transients; discard warm-up and report
sensitivity. (ii) \emph{Controllability is intervention-specific}: coordinated participant behavior,
not exogenous price moves, drives realistic stress in this generator. (iii) \emph{Tail fidelity is
transient}: validate generated tails against real stationarity before relying on them. EcoMD is a
useful, controllable, differentiable scenario generator; these results are the conditions under which
its output can be trusted. \textbf{Limitations.} We demonstrate the pitfall in EcoMD; we conjecture it
affects any rollout-scored generator initialized away from its stationary distribution---neural SDEs
and agent-based models alike---and recommend the warm-up audit whenever rollouts are scored, but
verifying its prevalence across model families is future work. The fidelity test is likewise limited
to free intraday (crypto) data.

{\small
\bibliographystyle{plainnat}
\bibliography{references}
}

\end{document}
